\chapter[Band 28: Halbgruppen]{Band 28 auf einen Blick}
\noindent\textbf{Halbgruppen}\par\medskip
Grundlage sind binäre Operationen sowie die Wörter und Klammerungsbäume
aus Band 27. In dieser Übersicht bedeutet \(A\cong B\) stets
Isomorphie der jeweils angegebenen Halbgruppen.
\(\mathcal P_+(A)=\powerset(A)\setminus\{\varnothing\}\) trägt das Mengenprodukt.
Ein neutrales Element wird nicht vorausgesetzt; Potenzen haben positive Exponenten.

\begin{BandUebersicht}
\textbf{Assoziative binäre Operation}
\UebFormel{\begin{aligned}
&\star:A\times A\to A,\\
&a,b,c\in A\ \vdash\
 (a\star b)\star c=a\star(b\star c).
\end{aligned}}Diese Axiome definieren \(\SemigroupAlg A\star\).
Die Linksfaltung eines Nichtleerwortes heißt \(\WordFold\star w\).
\UebQuellen{\FormulaRefAuto{SemigroupDef}[def];
\FormulaRefAuto{LeftWordFoldDef}[def]}
&
\textbf{Blöcke und Klammerungen}
Bei \(\SemigroupAlg A\star\) gilt:\UebFormel{\begin{aligned}
&u,v\in\NonemptyWords A\ \vdash\\
&\WordFold\star{\WordConcat u v}
 =\WordFold\star u\star\WordFold\star v,\\
&w\in\NonemptyWords A,\ S,T\in\WordBracketings A w\
 \vdash\\
&\qquad\TreeEvaluation\star S=\TreeEvaluation\star T.
\end{aligned}}Addition und Multiplikation auf \(\mathbb N\) sind Beispiele.
\UebQuellen{\FormulaRefAuto{SemigroupWordBlockLaw}[thm];
\FormulaRefAuto{SemigroupBracketingIndependence}[thm];
\FormulaRefAuto{NaturalAdditionSemigroup}[thm];
\FormulaRefAuto{NaturalMultiplicationSemigroup}[thm]} \\
\addlinespace[.8ex]

\textbf{Positive Potenzen und Idempotente}
Für \(\SemigroupAlg A\star\) und \(a\in A\):\UebFormel{\begin{aligned}
&a^1=a,\qquad a^{n+1}=a^n\star a\quad(n\ge1),\\
&E(A)=\{e\in A\mid e\star e=e\}.
\end{aligned}}
\UebQuellen{\FormulaRefAuto{SemigroupPositivePowerDef}[def];
\FormulaRefAuto{SemigroupPowerOne}[thm];
\FormulaRefAuto{SemigroupPowerSuccessor}[thm];
\FormulaRefAuto{SemigroupIdempotentSetDef}[def]}
&
\textbf{Potenzen rechnen assoziativ}
Für \(m,n\in\mathbb N\setminus\{0\}\):\UebFormel{\begin{aligned}
&\SemigroupAlg A\star,\ a\in A\ \vdash\\
&a^{m+n}=a^m\star a^n,\qquad
 a^m\star a^n=a^n\star a^m,\\
&\SemigroupAlg A\star,\ e\in E(A)\ \vdash\ e^n=e.
\end{aligned}}
\UebQuellen{\FormulaRefAuto{SemigroupPowerAdditionLaw}[thm];
\FormulaRefAuto{SemigroupPowersCommute}[thm];
\FormulaRefAuto{SemigroupIdempotentPower}[thm]} \\
\addlinespace[.8ex]

\textbf{Unterhalbgruppen und Erzeugung}
Bei \(\SemigroupAlg A\star\) heißt \(U\) Unterhalbgruppe, wenn\UebFormel{\begin{aligned}
&\varnothing\ne U\subseteq A,\quad
 \forall x,y\in U\ (x\star y\in U).\\
&\varnothing\ne X\subseteq A\ \vdash\\
&\langle X\rangle=\bigcap\mathcal U_X,\\
&\mathcal U_X=\{U\mid\mathsf{UHGr}(U;A,\star),\\
&\hspace{29mm}X\subseteq U\}.
\end{aligned}}
\UebQuellen{\FormulaRefAuto{SemigroupSubsemigroupDef}[def];
\FormulaRefAuto{SemigroupGeneratedSubsemigroupDef}[def]}
&
\textbf{Kleinste produktabgeschlossene Erweiterung}
Für die links angegebenen Voraussetzungen:\UebFormel{\begin{aligned}
&\mathsf{UHGr}(\langle X\rangle;A,\star),\
 X\subseteq\langle X\rangle,\\
&\mathsf{UHGr}(U;A,\star),\ X\subseteq U
 \ \vdash\ \langle X\rangle\subseteq U,\\
&a\in A\ \vdash\
 \langle a\rangle=\{a^n\mid n\in\mathbb N,\ n\ge1\}.
\end{aligned}}
\UebQuellen{\FormulaRefAuto{GeneratedSubsemigroupIsSubsemigroup}[thm];
\FormulaRefAuto{GeneratedSubsemigroupContainsGenerators}[thm];
\FormulaRefAuto{GeneratedSubsemigroupMinimal}[thm];
\FormulaRefAuto{SemigroupMonogenicPowerDescription}[thm]} \\
\addlinespace[.8ex]

\textbf{Ideale und zweiseitige Teilbarkeit}
In \(\SemigroupAlg A\star\) ist ein Ideal eine nichtleere
Teilmenge \(I\subseteq A\) mit \(AI\cup IA\subseteq I\).
Für \(a\in A\) setzt man\UebFormel{\begin{aligned}
&J(a)=\{a\}\cup Aa\cup aA\cup AaA,\\
&a\mid_A b\ \leftrightarrow\\
&\quad b=a\lor\exists x\in A\ (b=x\star a)\\
&\quad{}\lor\exists y\in A\ (b=a\star y)\\
&\quad{}\lor\exists x,y\in A\ (b=(x\star a)\star y).
\end{aligned}}
\UebQuellen{\FormulaRefAuto{SemigroupIdealDef}[def];
\FormulaRefAuto{SemigroupPrincipalTwoSidedIdealDef}[def];
\FormulaRefAuto{SemigroupTwoSidedDividesDef}[def]}
&
\textbf{Hauptideal und Teilbarkeit}
Für \(\SemigroupAlg A\star\) und \(a,b,c\in A\):\UebFormel{\begin{aligned}
&\mathsf{Id}(J(a);A,\star),\\
&a\mid_A b\ \leftrightarrow\ b\in J(a),\\
&a\mid_A a,\qquad
 (a\mid_A b\land b\mid_A c)\to a\mid_A c.
\end{aligned}}
\UebQuellen{\FormulaRefAuto{SemigroupPrincipalTwoSidedIdealIsIdeal}[thm];
\FormulaRefAuto{SemigroupTwoSidedDividesPrincipalIdeal}[thm];
\FormulaRefAuto{SemigroupDivisibilityReflexive}[thm];
\FormulaRefAuto{SemigroupTwoSidedDividesTransitive}[thm]} \\
\addlinespace[.8ex]

\textbf{Green-Relationen}
In einer festen Halbgruppe und für \(a,b\in A\) vergleichen
\(L(a),R(a),J(a)\) die linken, rechten und zweiseitigen Hauptideale:\UebFormel{\begin{aligned}
&a\mathcal L b\leftrightarrow L(a)=L(b),\\
&a\mathcal R b\leftrightarrow R(a)=R(b),\\
&a\mathcal J b\leftrightarrow J(a)=J(b),\\
&\mathcal H=\mathcal L\cap\mathcal R.
\end{aligned}}
\UebQuellen{\FormulaRefAuto{SemigroupGreenLDef}[def];
\FormulaRefAuto{SemigroupGreenRDef}[def];
\FormulaRefAuto{SemigroupGreenJDef}[def];
\FormulaRefAuto{SemigroupGreenHDef}[def]}
&
\textbf{Vier Äquivalenzrelationen}
Bei \(\SemigroupAlg A\star\) sind
\(\mathcal L,\mathcal R,\mathcal J,\mathcal H\) reflexiv, symmetrisch und
transitiv auf \(A\). Insbesondere:\UebFormel{\begin{aligned}
&a,b\in A\ \vdash\\
&a\mathcal J b\ \leftrightarrow\
 (a\mid_A b\land b\mid_A a).
\end{aligned}}
\UebQuellen{\FormulaRefAuto{SemigroupGreenLIsEquivalence}[thm];
\FormulaRefAuto{SemigroupGreenRIsEquivalence}[thm];
\FormulaRefAuto{SemigroupGreenJIsEquivalence}[thm];
\FormulaRefAuto{SemigroupGreenHIsEquivalence}[thm];
\FormulaRefAuto{SemigroupGreenJMutualTwoSidedDivisibility}[thm]} \\
\addlinespace[.8ex]

\textbf{Homomorphismen}
Für Halbgruppen \(A,B\):\UebFormel{\begin{aligned}
&f:A\to B,\\
&x,y\in A\ \vdash\\
&\quad f(x\star y)=f(x)\diamond f(y).
\end{aligned}}Ein Isomorphismus ist zusätzlich bijektiv.
\UebQuellen{\FormulaRefAuto{SemigroupHomDef}[def];
\FormulaRefAuto{SemigroupIsoDef}[def]}
&
\textbf{Komposition und Mengenbild}
Homomorphismen \(f:A\to B\), \(g:B\to C\) sind komponierbar.
Ein Isomorphismus \(f\) induziert durch \(X\mapsto f[X]\):\UebFormel{\begin{aligned}
&A\cong B\ \vdash\
 \mathcal P_+(A)\cong\mathcal P_+(B).
\end{aligned}}
\UebQuellen{\FormulaRefAuto{SemigroupHomComposition}[thm];
\FormulaRefAuto{SemigroupIsoInducesPowerSemigroupIso}[thm]} \\
\addlinespace[.8ex]

\textbf{Potenzhalbgruppe und Einermengenkopie}
Für \(\SemigroupAlg A\star\) und \(X,Y\in\mathcal P_+(A)\):
\UebFormel{\begin{aligned}
&X\star_{\mathcal P}Y
 =\{x\star y\mid x\in X,\ y\in Y\},\\
&\mathcal S(A)=\{\{a\}\mid a\in A\}.
\end{aligned}}
\UebQuellen{\FormulaRefAuto{PowerSemigroupProductDef}[def];
\FormulaRefAuto{SemigroupSingletonCopyDef}[def]}
&
\textbf{Mengenprodukte bleiben assoziativ}
\UebFormel{\begin{aligned}
&\SemigroupAlg A\star\ \vdash\\
&\quad\SemigroupAlg{\mathcal P_+(A)}{\star_{\mathcal P}},\\
&A\cong\mathcal S(A),\qquad
 \{a\}\star_{\mathcal P}\{b\}=\{a\star b\}.
\end{aligned}}Die letzte Gleichung gilt für \(a,b\in A\).
\UebQuellen{\FormulaRefAuto{NonemptyPowerSemigroup}[thm];
\FormulaRefAuto{SemigroupSingletonCopyIsomorphism}[thm];
\FormulaRefAuto{PowerSemigroupProductMembership}[thm]} \\
\addlinespace[.8ex]

\textbf{Koordinaten an einem Basiselement}
Für \(\SemigroupAlg A\star\), \(e\in A\):\UebFormel{\begin{aligned}
&\mathcal K_e=(Ae)\times(eA),\\
&(p,q)\square_e(r,s)=(p,s),\\
&\chi_e(x)=(x\star e,e\star x).
\end{aligned}}
\UebQuellen{\FormulaRefAuto{SemigroupCoordinateCarrierDef}[def];
\FormulaRefAuto{SemigroupCoordinateProductDef}[def];
\FormulaRefAuto{SemigroupCoordinateMapDef}[def]}
&
\textbf{Kriterium für eine Koordinatenisomorphie}
Unter den linken Voraussetzungen genügen\UebFormel{\begin{aligned}
&\forall x\in A\ (x\star x=x),\\
&\forall x,y,z\in A\ ((x\star y)\star z=x\star z)\\
&\qquad\vdash\ \SemigroupCoordinateIso A\star e.
\end{aligned}}
\UebQuellen{\FormulaRefAuto{SemigroupCoordinateIsoCriterion}[thm]} \\
\addlinespace[.8ex]

\textbf{Inflation}
Eine Halbgruppe \(S\supseteq G\) ist eine Inflation von \(G\),
wenn ein Rückzug \(r:S\to G\) erfüllt:\UebFormel{\begin{aligned}
&r|_G=\operatorname{id}_G,\\
&x,y\in S\ \vdash\ x\star y=r(x)\star r(y).
\end{aligned}}
\UebQuellen{\FormulaRefAuto{SemigroupInflationDef}[def]}
&
\textbf{Rekonstruktion durch passende Fasern}
Für Inflationen \(r:S\to G\), \(s:T\to H\) und
einen Halbgruppenisomorphismus \(\varphi:G\to H\) gilt:\UebFormel{\begin{aligned}
&\forall g\in G\quad
 r^{-1}[\{g\}]\sim s^{-1}[\{\varphi(g)\}]\\
&\qquad\vdash\ S\cong T.
\end{aligned}}\(\sim\) bedeutet hier Existenz einer Bijektion.
\UebQuellen{\FormulaRefAuto{InflationFiberIsomorphism}[thm]} \\
\addlinespace[.8ex]

\textbf{Das Mogiljanskaja-Paar}
Für die disjunkten Schichten \(L=\{a_i\}\), \(D=\{b_{ij}\}\)
aus Band 21:\UebFormel{\begin{aligned}
&S=L\cup D\cup\{\mathbf0\},\quad T=S\cup\{k\},\\
&a_i a_j=b_{ij},\qquad
 xy=\mathbf0\ \text{sonst}.
\end{aligned}}Die markierten Elemente \(\mathbf0,k\) liegen außerhalb der Schichten.
Beide Operationen folgen derselben Regel.
\UebQuellen{\FormulaRefAuto{MogiljanskajaPairDef}[def];
\FormulaRefAuto{MogiljanskajaRuleDef}[def]}
&
\textbf{Unendliches Gegenbeispiel, einschließlich formalem Anhang}
\UebFormel{\begin{aligned}
&\neg\Finite S,\qquad\neg\Finite T,\\
&S\not\cong T,\qquad
 \mathcal P_+(S)\cong\mathcal P_+(T).
\end{aligned}}Die im Anhang ausgeführte Bijektion \(\Phi\) erhält sämtliche
Mengenprodukte. Dies widerlegt die uneingeschränkte Rekonstruktion.
\UebQuellen{\FormulaRefAuto{MogiljanskajaPairCounterexample}[thm];
\FormulaRefAuto{MogiljanskajaArgumentM16PowerIsomorphism}[thm];
\FormulaRefAuto{MogiljanskajaArgumentM17CounterexampleWitness}[thm]} \\
\addlinespace[.8ex]
\end{BandUebersicht}
